\documentclass[12pt,a4paper]{article}

% --- Page Setup & Margins ---
\usepackage[
  a4paper,
  left=25mm,
  right=25mm,
  top=20mm,
  bottom=20mm
]{geometry}

% --- Core Packages ---
\usepackage{graphicx}
\usepackage{amsmath,amssymb}
\usepackage{array,tabularx}
\usepackage{enumitem}
\usepackage{enumitem}
\usepackage{calc}

% --- Typography ---
\usepackage{helvet}
\renewcommand{\familydefault}{\sfdefault} % Clean Sans-Serif font

% --- Disable Headers & Footers Completely ---
\pagestyle{empty}

\setlength{\parindent}{0pt}
\setlength{\parskip}{4pt}

% --- Custom Exam Macros ---
\newcommand{\qnum}[1]{\vspace{6pt}\noindent\textbf{#1}\hspace{0.5em}}
\newcommand{\qmarks}[1]{\hfill\textbf{[#1]}}
\newcommand{\totalmarks}[1]{\par\nopagebreak\hfill\textbf{[Total: #1]}\par}

% Dotted line for text responses
\newcommand{\wline}{\par\vspace{1.5pt}\noindent\dotfill\par\vspace{1.5pt}}

% Answer field with label, dots, unit, and marks at right edge
\newcommand{\ansfield}[3]{%
  \par\vspace{6pt}
  \noindent #1 = \dotfill\ \mathrm{#2} \quad [#3]\par\vspace{4pt}
}

\begin{document}
\begin{center}
	\textsc{This document contains Topic list and the Question paper}\\
	\textsc{*Each question might contain more than one topic }
\end{center}

\textbf{Topic List}
\begin{enumerate}[label=\textbf{Q\arabic*:}]
    \item Motion or Kinematics, Forces or Dynamics
    \item Forces or Dynamics, Work Energy and Power
    \item Forces or Dynamics, Mass Weight and Density, Momentum
    \item Kinetic Model of Matter, Thermal Properties \& Temperature
    \item Electromagnetic Spectrum
    \item Electric Charge, Electrical Quantities
    \item Electrical Quantities, Practical Electricity / Electrical Safety
    \item The Nuclear Atom, Radioactivity
    \item The Stars \& Universe
\end{enumerate}
\newpage

% =========================================================
% PAGE 2 -- Question 1
% =========================================================

\qnum{1} A car travels along a straight horizontal road, initially at a constant speed. The speed--time graph in Fig.~1.1 shows the motion of the car from time $t = 0$.

\begin{center}
  \includegraphics[width=0.90\linewidth]{q1_fig1.1} \\ \vspace{4pt}
  \textbf{Fig. 1.1}
\end{center}

At time $t = 5.7\,\mathrm{s}$, the driver sees an obstruction in the road.

The driver applies the brakes a short time later, and the car begins to slow down.

\textbf{(a)} Using Fig.~1.1, determine:

\hspace*{1.5em}\textbf{(i)} the speed of the car at time $t = 0$
\ansfield{speed}{m/s}{1}

\hspace*{1.5em}\textbf{(ii)} the distance travelled by the car between the driver seeing the obstruction at $t = 5.7\,\mathrm{s}$ and the car beginning to slow down.

\vspace{2cm}
\ansfield{distance}{m}{2}

\newpage

% =========================================================
% PAGE 3 -- Question 1 (Continued)
% =========================================================

\hspace*{1.5em}\textbf{(iii)} the distance travelled by the car between $t = 0$ and $t = 11.2\,\mathrm{s}$.

\vspace{4cm}
\ansfield{distance}{m}{3}

\textbf{(b)} The speed of the car affects both the thinking distance and the braking distance.

State \textbf{one} factor that affects \textbf{only} the braking distance of the car.

Explain how a change in this factor increases the braking distance.

\vspace{4pt}
\noindent factor \dotfill \wline
\noindent explanation \dotfill \wline \wline
\hfill[3]

\totalmarks{9}

\newpage

% =========================================================
% PAGE 4 -- Question 2
% =========================================================

\qnum{2} A spherical container is used to carry measuring instruments to the bottom of a lake.

Fig.~2.1 shows the container held at rest just below the water surface of a lake.

\begin{center}
  \includegraphics[width=0.80\linewidth]{q2_fig2.1} \\ \vspace{4pt}
  \textbf{Fig. 2.1} (not to scale)
\end{center}

The container is released and falls through the water to the bottom of the lake.

\textbf{(a)} The container reaches terminal velocity before it hits the bottom of the lake.

Explain, in terms of the forces on the container, the motion of the container through the water until it reaches terminal velocity.

\wline \wline \wline \wline \wline
\hfill[3]

\newpage

% =========================================================
% PAGE 5 -- Question 2 (Continued)
% =========================================================

\textbf{(b)} The container falls through a vertical distance of $32\,\mathrm{m}$ before it hits the bottom of the lake.

The container and instruments have a total mass of $74\,\mathrm{kg}$.

\hspace*{1.5em}\textbf{(i)} Calculate the energy transferred from the gravitational potential energy store of the container.

\vspace{2.5cm}
\ansfield{energy transferred}{J}{2}

\hspace*{1.5em}\textbf{(ii)} State the principle of conservation of energy.

\wline \wline \wline
\hfill[2]

\hspace*{1.5em}\textbf{(iii)} Work is done to transfer energy from the gravitational potential energy store as the container falls at terminal velocity.

Describe, in terms of the work done, how the energy is transferred from the gravitational potential energy store to \textbf{one} other energy store in the water.

\wline \wline \wline
\hfill[2]

\totalmarks{9}

\newpage

% =========================================================
% PAGE 6 -- Question 3
% =========================================================

\qnum{3} Fig.~3.1 shows a firefighter directing a jet of water at a wall.

\begin{center}
  \includegraphics[width=0.8\linewidth]{q3_fig3.1} \\ \vspace{4pt}
  \textbf{Fig. 3.1}
\end{center}

In $1.0\,\mathrm{s}$, a volume of $7.5 \times 10^{-3}\,\mathrm{m}^3$ of water hits the wall horizontally, at a speed of $24\,\mathrm{m/s}$.

The density of water is $1000\,\mathrm{kg/m}^3$.

\textbf{(a)} Calculate:

\hspace*{1.5em}\textbf{(i)} the mass of the water that hits the wall in $1.0\,\mathrm{s}$
\vspace{2cm}
\ansfield{mass}{kg}{2}

\hspace*{1.5em}\textbf{(ii)} the horizontal momentum of the water that hits the wall in $1.0\,\mathrm{s}$.
\vspace{2cm}
\ansfield{momentum}{kg\,m/s}{2}

\textbf{(b)} The horizontal momentum of the water decreases to zero when the water hits the wall. None of the water bounces back from the wall.

Explain why the momentum calculated in \textbf{(a)(ii)} is equal in size to the horizontal force exerted on the water by the wall.

\wline \wline
\hfill[2]

\newpage

% =========================================================
% PAGE 7 -- Question 3 (Continued)
% =========================================================

\textbf{(c)} Explain, in terms of Newton's third law, why there is a horizontal force exerted on the wall by the water.

\wline \wline \wline
\hfill[2]

\textbf{(d)} The pump that forces the water to flow through the hose is adjusted, and the speed of the water leaving the hose doubles.

Explain why the force exerted by the water on the wall increases by a factor of 4.

\wline \wline
\hfill[1]

\totalmarks{9}

\newpage

% =========================================================
% PAGE 8 -- Question 4
% =========================================================

\qnum{4} A large test-tube contains a thermometer and some solid wax at a temperature of $21\,^{\circ}\mathrm{C}$.

At time $t = 0$, the test-tube is partially immersed in a beaker of boiling water and the temperature of the wax is recorded for the next $600\,\mathrm{s}$.

Fig.~4.1 shows the variation of the temperature of the wax with time $t$.

\begin{center}
  \includegraphics[width=0.99\linewidth]{q4_fig4.1} \\ \vspace{4pt}
  \textbf{Fig. 4.1}
\end{center}

\textbf{(a)} Explain how Fig.~4.1 shows that the melting temperature of the wax is $63\,^{\circ}\mathrm{C}$.

\wline \wline
\hfill[1]

\newpage

% =========================================================
% PAGE 9 -- Question 4 (Continued)
% =========================================================

\textbf{(b)} As the temperature of the wax increases, the energy in the internal energy store increases.

\hspace*{1.5em}\textbf{(i)} The mass of the wax in the test-tube is $0.040\,\mathrm{kg}$ and its temperature at time $t = 0$ is $21\,^{\circ}\mathrm{C}$. The specific heat capacity of wax is $2100\,\mathrm{J/(kg\,^{\circ}C)}$.

Using Fig.~4.1, determine the increase in the energy in the internal energy store of the wax between $t = 0$ and $t = 220\,\mathrm{s}$. Show your working.

\vspace{3cm}
\ansfield{increase in energy in internal energy store}{J}{3}

\hspace*{1.5em}\textbf{(ii)} Describe what happens to the motion of the molecules of the solid wax as the internal energy of the wax increases.

\wline \wline \wline
\hfill[2]

\textbf{(c)} When the wax reaches its melting temperature, energy continues to be transferred thermally to the wax.

Describe, in terms of the forces between particles, why energy is required to melt the wax.

\wline \wline \wline
\hfill[2]

\totalmarks{8}

\newpage

% =========================================================
% PAGE 10 -- Question 5
% =========================================================

\qnum{5} The microwave region is one region of the electromagnetic spectrum.

\textbf{(a)} Fig.~5.1 represents the electromagnetic spectrum divided into the seven main regions.

\begin{center}
  \includegraphics[width=0.90\linewidth]{q5_fig5.1} \\ \vspace{4pt}
  \textbf{Fig. 5.1}
\end{center}

The regions are arranged from left to right in order of increasing wavelength.

On Fig.~5.1:
\begin{itemize}[leftmargin=2.5em,itemsep=1pt]
  \item indicate the microwave region by marking it with the letter \textbf{M}
  \item indicate the ultraviolet region by marking it with the letters \textbf{UV}.
\end{itemize}
\hfill[1]

\textbf{(b)} The frequency of the microwaves used by some satellite television systems is $12\,\mathrm{GHz}$.

\hspace*{1.5em}\textbf{(i)} The speed of microwaves in a vacuum is $3.0 \times 10^{8}\,\mathrm{m/s}$.

Calculate the wavelength of these microwaves in a vacuum.

\vspace{2.5cm}
\ansfield{wavelength}{m}{3}

\hspace*{1.5em}\textbf{(ii)} Describe how microwaves are used in satellite television systems.

\wline \wline \wline
\hfill[2]

\newpage

% =========================================================
% PAGE 11 -- Question 5 (Continued)
% =========================================================

\textbf{(c)} All microwave ovens have a switch that turns the oven off when the door is opened.

Suggest why it is important for a microwave oven to switch off when the door is opened.

\wline \wline \wline
\hfill[2]

\totalmarks{8}

\newpage

% =========================================================
% PAGE 12 -- Question 6
% =========================================================

\qnum{6} An electric car is driven by a direct current (d.c.) motor that is powered by a direct current power supply.

\textbf{(a)} Describe how a direct current differs from an alternating current (a.c.).

\wline \wline
\hfill[1]

\textbf{(b)} Fig.~6.1 represents the motor circuit. Point P is a point in the circuit between the motor and the positive terminal of the power supply.

Point P is shown to the left of the motor in the circuit diagram.

\begin{center}
  \includegraphics[width=0.55\linewidth]{q6_fig6.1} \\ \vspace{4pt}
  \textbf{Fig. 6.1}
\end{center}

During a time of $120\,\mathrm{s}$, the current in the circuit at point P is $41\,\mathrm{A}$.

\hspace*{1.5em}\textbf{(i)} Calculate the charge that flows through point P in the time of $120\,\mathrm{s}$.

\vspace{2.5cm}
\ansfield{charge}{C}{2}

\newpage

% =========================================================
% PAGE 13 -- Question 6 (Continued)
% =========================================================

\hspace*{1.5em}\textbf{(ii)} The current in the circuit is due to the flow of electrons. Each electron carries a charge of magnitude $1.6 \times 10^{-19}\,\mathrm{C}$.

State the direction which the electrons move past point P and determine the number of electrons that pass point P in $1.0\,\mathrm{s}$.

\vspace{4pt}
\noindent direction \dotfill

\vspace{2.5cm}
\noindent number of electrons = \dotfill \hfill[2]

\textbf{(c)} The power supply for the electric car is a number of batteries, each made from a large number of identical cells.

The e.m.f. (electromotive force) of each cell is $3.7\,\mathrm{V}$.

\hspace*{1.5em}\textbf{(i)} Define the term \textit{electromotive force}.

\wline \wline
\hfill[2]

\hspace*{1.5em}\textbf{(ii)} Each battery consists of 92 cells in series.

Calculate the e.m.f. of one battery.

\vspace{1.5cm}
\ansfield{e.m.f.}{V}{2}

\hspace*{1.5em}\textbf{(iii)} The power supply is made by connecting 85 of these batteries in parallel.

State the e.m.f. of the power supply.

\vspace{0.8cm}
\ansfield{e.m.f.}{V}{1}

\totalmarks{10}

\newpage

% =========================================================
% PAGE 14 -- Question 7
% =========================================================

\qnum{7} An electric kettle has an outer casing made of metal. Fig.~7.1 shows the kettle.

\begin{center}
  \includegraphics[width=0.80\linewidth]{q7_fig7.1} \\ \vspace{4pt}
  \textbf{Fig. 7.1}
\end{center}

The heater in the electric kettle is connected to a $230\,\mathrm{V}$ mains power supply and has a power rating of $2300\,\mathrm{W}$.

\textbf{(a)} The kettle is switched on.

\hspace*{1.5em}\textbf{(i)} Calculate the current in the heater.

\vspace{2.5cm}
\ansfield{current}{A}{2}

\hspace*{1.5em}\textbf{(ii)} The manufacturer of the kettle states that the cable for the kettle is safe for currents that are smaller than $15\,\mathrm{A}$. The wiring in the mains power supply in the wall is safe for currents that are smaller than $20\,\mathrm{A}$.

The fuses available for the kettle have the ratings shown.

\begin{center}
\textbf{3A \quad 5A \quad 7A \quad 10A \quad 13A \quad 15A \quad 18A \quad 20A \quad 25A}
\end{center}

State which of these fuses is the most appropriate. Explain why it is the most appropriate.

\vspace{4pt}
\noindent most appropriate fuse rating = \dotfill\ \mathrm{A}

\noindent explanation \dotfill \wline \wline
\hfill[2]

\newpage

% =========================================================
% PAGE 15 -- Question 7 (Continued)
% =========================================================

\textbf{(b)} There are three wires in the cable that connects the kettle to the mains power supply:
\begin{itemize}[leftmargin=2.5em,itemsep=1pt]
  \item the earth wire
  \item the live (line) wire
  \item the neutral wire.
\end{itemize}

The wires are all correctly connected.

\hspace*{1.5em}\textbf{(i)} State the name of the wire that the fuse is connected into. Explain why the fuse is connected into this wire.

\vspace{4pt}
\noindent name of wire \dotfill

\noindent explanation \dotfill \wline \wline
\hfill[2]

\hspace*{1.5em}\textbf{(ii)} The insulation on the live wire is damaged so that the wire is exposed.

Explain what happens when the exposed live wire touches the metal outer casing.

\wline \wline \wline
\hfill[2]

\textbf{(c)} The kettle is switched on for a total time of $34\,\mathrm{h}$ during one year.

The cost of electricity is \$0.32 per kilowatt-hour (kWh).

Calculate the cost of using the kettle during the year.

\vspace{2.5cm}
\noindent $\mathrm{cost} = \$\dotfill$ \hfill[2]

\totalmarks{9}

\newpage

% =========================================================
% PAGE 16 -- Question 8
% =========================================================

\qnum{8} The element lithium has several different isotopes.

\textbf{(a)} \hspace*{0.3em}\textbf{(i)} State \textbf{one} way in which the compositions of the atoms of all lithium isotopes are the same.

\wline \wline
\hfill[1]

\hspace*{1.5em}\textbf{(ii)} State \textbf{one} way in which the composition of a lithium atom of one isotope is different from the composition of a lithium atom of a different isotope.

\wline \wline
\hfill[1]

\textbf{(b)} Fig.~8.1 is a diagram that represents a neutral atom of a radioactive isotope of lithium ($\mathrm{Li}$).

\begin{center}
  \includegraphics[width=0.7\linewidth]{q8_fig8.1} \\ \vspace{4pt}
  \textbf{Fig. 8.1}
\end{center}

There are three electrons in orbit around the nucleus.

\hspace*{1.5em}\textbf{(i)} The nuclide symbols for isotopes of lithium are represented as $^{X}_{Y}\mathrm{Li}$ where $X$ and $Y$ are numbers.

Determine the values of $X$ and $Y$ for the isotope represented by Fig.~8.1.

\vspace{6pt}
\noindent $X = $ \dotfill

\noindent $Y = $ \dotfill
\hfill[2]

\hspace*{1.5em}\textbf{(ii)} An ion of this isotope of lithium has a single positive charge and is represented as $\left(^{X}_{Y}\mathrm{Li}^{+}\right)$.

Describe how the diagram that represents this ion differs from Fig.~8.1.

\wline \wline
\hfill[1]

\newpage

% =========================================================
% PAGE 17 -- Question 8 (Continued)
% =========================================================

\textbf{(c)} Radioactive isotopes emit nuclear radiation which can affect living things.

\hspace*{1.5em}\textbf{(i)} State \textbf{one} damaging effect of nuclear radiation.

\wline \wline
\hfill[1]

\hspace*{1.5em}\textbf{(ii)} The emission of radioactive radiation from an unstable isotope is described as random and spontaneous.

State what is meant by `random' and what is meant by `spontaneous'.

\vspace{4pt}
\noindent random \dotfill \wline

\noindent spontaneous \dotfill \wline
\hfill[2]

\totalmarks{8}

\newpage

% =========================================================
% PAGE 18 -- Question 9
% =========================================================

\qnum{9} Fig.~9.1 shows how the Sun and the rest of the Solar System orbit around the centre of the Milky Way in a circular path.

\begin{center}
  \includegraphics[width=0.8\linewidth]{q9_fig9.1} \\ \vspace{4pt}
  \textbf{Fig. 9.1}
\end{center}

The Sun travels around the circular path at a speed $v$.

\textbf{(a)} The speed of light in a vacuum is $9.5 \times 10^{12}\,\mathrm{km/year}$.

The Sun is $26\,000$ light-years from the centre of the Milky Way.

\hspace*{1.5em}\textbf{(i)} Determine the distance from the Sun to the centre of the Milky Way in kilometres ($\mathrm{km}$).

\vspace{2.5cm}
\ansfield{distance}{km}{2}

\hspace*{1.5em}\textbf{(ii)} It takes the Sun $2.3 \times 10^{8}$ years to complete one orbit around the centre of the Milky Way.

Calculate the speed $v$.

\vspace{2.5cm}
\ansfield{$v$}{m/s}{2}

\newpage

% =========================================================
% PAGE 19 -- Question 9 (Continued)
% =========================================================

\textbf{(b)} Astronomical observations suggest that, at the centre of the Milky Way, there is a black hole.

Black holes are produced from red supergiants.

\hspace*{1.5em}\textbf{(i)} Near the end of its life, a massive star becomes a red supergiant.

Describe what happens inside the star as it becomes a red supergiant.

\wline \wline \wline
\hfill[2]

\hspace*{1.5em}\textbf{(ii)} Describe how a red supergiant produces a black hole.

\wline \wline \wline \wline
\hfill[3]

\hspace*{1.5em}\textbf{(iii)} Describe how the heaviest elements are produced.

\wline \wline
\hfill[1]

\totalmarks{10}

\end{document}